\%============================================================
\chapter{Bài Học: Khẳng Định Trong Phủ Định, Phủ Định Trong Khẳng Định (Affirmation in Negation, Negation in Affirmation)}
\begin{center}
  \textit{\color{truyengold}``Rejecting what is wrong is the first step in affirming what is right.''}\\[4pt]
  \textit{(Từ chối cái sai là bước đầu tiên để khẳng định cái đúng.)}\\[4pt]
  \small\color{truyenblue}--- Cổ Kim Đông Tây ---
\end{center}
\ngancach

\section{Lời Khẳng Định Bằng Sự Im Lặng}
\begin{truyen}{Lời Khẳng Định Bằng Sự Im Lặng}{Thomas More --- Anh, 1535}
\chuhoa{K}{hi} Vua Henry VIII tự phong làm người đứng đầu Giáo hội Anh, ngài yêu cầu toàn bộ quần thần phải tuyên thệ ủng hộ.\\[4pt]
\textit{(When King Henry VIII declared himself Supreme Head of the Church of England, he required all his subjects to swear an oath of support.)}

Thomas More, cựu Tể tướng, đã từ chối tuyên thệ, nhưng ông cũng không mở miệng chỉ trích nhà vua.\\[4pt]
\textit{(Thomas More, the former Lord Chancellor, refused to take the oath, but he also kept his mouth shut from criticizing the king.)}

Ông chọn sự im lặng tuyệt đối; nhà vua cho rằng sự im lặng đó là một sự phủ định đầy khiêu khích.\\[4pt]
\textit{(He chose absolute silence; the king considered that silence to be a provocative negation.)}

Thomas More bị kết án tử hình vì sự im lặng của mình.\\[4pt]
\textit{(Thomas More was sentenced to death for his silence.)}

Nhưng bằng cách phủ định không nói ra lời thề dối trá, ông đã khẳng định mạnh mẽ nhất lương tâm và đức tin của mình trước toàn nước Anh.\\[4pt]
\textit{(But by negating to speak a false oath, he made the strongest affirmation of his conscience and faith before all of England.)}

\end{truyen}
\begin{baihoc}
  Đôi khi, từ chối tham gia vào sự dối trá chính là lời khẳng định đanh thép nhất cho sự thật.\\[4pt]
  \textit{(Sometimes, refusing to participate in a lie is the most resounding affirmation of the truth.)}
\end{baihoc}

\section{Chữ ``Không'' Của Rosa Parks}
\begin{truyen}{Chữ ``Không'' Của Rosa Parks}{Rosa Parks --- Hoa Kỳ, 1955}
\chuhoa{N}{ăm} 1955 tại Montgomery, luật pháp quy định người da đen phải nhường ghế cho người da trắng trên xe buýt.\\[4pt]
\textit{(In 1955 in Montgomery, the law required Black people to give up their bus seats to white passengers.)}

Một buổi chiều muộn, một tài xế yêu cầu Rosa Parks đứng lên nhường ghế, cô chỉ nói một từ: ``Không.''\\[4pt]
\textit{(One late afternoon, a driver ordered Rosa Parks to stand up and give her seat, she said only one word: `No.')}

Sự phủ định ngắn gọn của cô đã dẫn đến việc cô bị bắt giam ngay lập tức.\\[4pt]
\textit{(Her brief negation led to her immediate arrest and imprisonment.)}

Nhưng chính chữ ``Không'' đó đã châm ngòi cho cuộc Tẩy chay xe buýt Montgomery kéo dài 381 ngày.\\[4pt]
\textit{(But that very 'No' ignited the 381-day Montgomery Bus Boycott.)}

Bằng cách phủ định một đạo luật bất công, cô đã khẳng định quyền bình đẳng và nhân phẩm của hàng triệu người.\\[4pt]
\textit{(By negating an unjust law, she affirmed the equal rights and human dignity of millions.)}

\end{truyen}
\begin{baihoc}
  Một chữ ``Không'' đúng lúc có thể kiến tạo lại cả một xã hội.\\[4pt]
  \textit{(A well-timed 'No' can rebuild an entire society.)}
\end{baihoc}

\section{Steve Jobs Từ Chối Hàng Ngàn Yù Tưởng}
\begin{truyen}{Steve Jobs Từ Chối Hàng Ngàn Yù Tưởng}{Steve Jobs --- Hoa Kỳ, 1997}
\chuhoa{K}{hi} Steve Jobs trở lại Apple năm 1997, công ty đang sản xuất hàng chục dòng máy tính và máy in chồng chéo nhau.\\[4pt]
\textit{(When Steve Jobs returned to Apple in 1997, the company was producing dozens of overlapping computer and printer lines.)}

Ông tổ chức một cuộc họp, vẽ một lưới hai nhân hai lên bảng trắng và nói: ``Chúng ta chỉ làm bốn sản phẩm.''\\[4pt]
\textit{(He held a meeting, drew a two-by-two grid on the whiteboard and said: `We are only making four products.')}

Ông đã phủ định hơn 300 sản phẩm đang phát triển, cắt giảm 70 phần trăm danh mục của Apple.\\[4pt]
\textit{(He negated over 300 products in development, cutting 70 percent of Apple's portfolio.)}

Mọi người kinh ngạc, nhưng Jobs giải thích: ``Tập trung có nghĩa là nói `KHÔNG' với một ngàn ý tưởng tốt khác.''\\[4pt]
\textit{(People were shocked, but Jobs explained: `Focus means saying NO to a thousand other good ideas.')}

Sự phủ định tàn nhẫn ấy đã giúp Apple dồn toàn lực vào những sản phẩm thực sự vĩ đại, cứu công ty khỏi bờ vực phá sản.\\[4pt]
\textit{(That ruthless negation allowed Apple to put all its strength into truly great products, saving the company from the brink of bankruptcy.)}

\end{truyen}
\begin{baihoc}
  Để khẳng định một điều xuất sắc, bạn phải dũng cảm phủ định hàng trăm điều tốt.\\[4pt]
  \textit{(To affirm one excellent thing, you must bravely negate a hundred good things.)}
\end{baihoc}

\section{Lời Phủ Định Của Socrates}
\begin{truyen}{Lời Phủ Định Của Socrates}{Socrates --- Hy Lạp, 399 TCN}
\chuhoa{K}{hắp} thành Athens, những nhà triết học tự xưng là người thông thái và biết mọi thứ trên đời.\\[4pt]
\textit{(Throughout Athens, philosophers claimed to be wise and to know everything in the world.)}

Nhưng Socrates lại đi khắp nơi với một lời phủ định nổi tiếng: ``Tôi chỉ biết một điều, đó là tôi không biết gì cả.''\\[4pt]
\textit{(But Socrates walked everywhere with a famous negation: `I know only one thing, that I know nothing.')}

Vì ông cho rằng mình không biết gì, ông liên tục đặt câu hỏi, thách thức những định kiến đã được mặc định là đúng.\\[4pt]
\textit{(Because he assumed he knew nothing, he constantly asked questions, challenging assumptions taken for granted as true.)}

Lời sấm truyền ở đền Delphi đã gọi ông là người thông thái nhất Athens --- chính vì ông nhận bết được sự ngu dốt của mình.\\[4pt]
\textit{(The Oracle of Delphi named him the wisest man in Athens --- precisely because he recognized his own ignorance.)}

Sống trong tư thế từ chối ảo tưởng tri thức, ông đã khẳng định được bản chất đích thực của trí tuệ.\\[4pt]
\textit{(Living in a posture that rejected the illusion of knowledge, he affirmed the true nature of wisdom.)}

\end{truyen}
\begin{baihoc}
  Biết giới hạn của bản thân là bước đầu tiên của trí tuệ thực sự.\\[4pt]
  \textit{(Knowing one's own limits is the first step of true wisdom.)}
\end{baihoc}

\section{Tuyên Ngôn Độc Lập --- Phủ Định Ách Thống Trị}
\begin{truyen}{Tuyên Ngôn Độc Lập --- Phủ Định Ách Thống Trị}{Thomas Jefferson --- Hoa Kỳ, 1776}
\chuhoa{T}{rong} Tuyên ngôn Độc lập Mỹ, Thomas Jefferson đã bắt đầu bằng việc liệt kê hàng loạt các hành động bạo ngược của Vua Anh.\\[4pt]
\textit{(In the American Declaration of Independence, Thomas Jefferson began by listing the serial tyrannies of the British King.)}

Văn kiện này thực chất là một bản cáo trạng lịch sử, một sự phủ định hoàn toàn quyền lực của Đế quốc Anh đối với vùng thuộc địa.\\[4pt]
\textit{(The document was essentially a historical indictment, a complete negation of the British Empire's authority over the colonies.)}

Nhưng cũng chính vì phủ định chế độ quân chủ chuyên chế đó, văn kiện đã đưa ra một khẳng định vĩ đại chưa từng có.\\[4pt]
\textit{(But precisely because it negated that tyrannical monarchy, the document made an unprecedented and great affirmation.)}

Đó là: ``Mọi người sinh ra đều bình đẳng... được tạo hóa ban cho những quyền không thể xâm phạm, trong đó có quyền Sống, quyền Tự do và mưu cầu Hạnh phúc.''\\[4pt]
\textit{(That is: `All men are created equal... endowed by their Creator with certain unalienable Rights, that among these are Life, Liberty and the pursuit of Happiness.')}

Sự bác bỏ cái sai đã làm nền tảng vững chắc để xây lên nền dân chủ tự do.\\[4pt]
\textit{(The rejection of the wrong served as the solid foundation to build a free democracy.)}

\end{truyen}
\begin{baihoc}
  Không đập bỏ những xiềng xích mục nát, ta không thể xây dựng được đài tự do mới.\\[4pt]
  \textit{(Without breaking the rotting chains, we cannot build a new monument of freedom.)}
\end{baihoc}

\section{Galileo Và Chuyển Động Của Trái Đất}
\begin{truyen}{Galileo Và Chuyển Động Của Trái Đất}{Galileo Galilei --- Ý, 1633}
\chuhoa{T}{rong} nhiều thế kỷ, Giáo hội tin rằng Trái đất là trung tâm của vũ trụ và mọi vì sao phải xoay quanh nó.\\[4pt]
\textit{(For centuries, the Church believed the Earth was the center of the universe and all stars must revolve around it.)}

Galileo quan sát thiên văn và nhận ra sự thật; ông bị đưa ra tòa án dị giáo và bị đe dọa thiêu sống nếu không từ bỏ quan điểm của mình.\\[4pt]
\textit{(Galileo observed the heavens and discovered the truth; he was brought before the Inquisition and threatened with burning at the stake if he did not renounce his view.)}

Vì già yếu không muốn chịu nhục hình, ông đành phải đọc lời phủ định chính khám phá khoa học của mình trước công chúng.\\[4pt]
\textit{(Being old and unwilling to suffer torture, he had to publicly read a negation of his own scientific discovery.)}

Nhưng ngay lúc bước xuống bục, ông khẽ thì thầm một lời khẳng định bất hủ: ``Dù sao thì trái đất vẫn quay.''\\[4pt]
\textit{(But exactly as he stepped down the podium, he muttered an immortal affirmation: `And yet it moves.')}

Sự phủ định bằng môi miệng không thể xóa sổ chân lý khách quan mà ông đã khẳng định bằng lý trí.\\[4pt]
\textit{(A verbal negation could not erase the objective truth he had affirmed with reason.)}

\end{truyen}
\begin{baihoc}
  Sự thật không thay đổi chỉ vì con người bị ép buộc phải chối bỏ nó.\\[4pt]
  \textit{(Truth does not change just because people are forced to deny it.)}
\end{baihoc}

\section{Bức Tường Berlin Sụp Đổ}
\begin{truyen}{Bức Tường Berlin Sụp Đổ}{Đức, 1989}
\chuhoa{B}{ức} tường Berlin là biểu tượng vật lý lớn nhất của sự ngăn cách, một sự phủ định tàn nhẫn đối với quyền tự do đi lại và đoàn tụ của người dân Đức.\\[4pt]
\textit{(The Berlin Wall was the greatest physical symbol of separation, a cruel negation of the German people's right to move freely and reunite.)}

Trong 28 năm, nó đứng đó từ chối khát vọng tự do, những người cố trèo qua đều bị bắn chết.\\[4pt]
\textit{(For 28 years, it stood there denying the desire for freedom, those who tried to climb over were shot dead.)}

Nhưng chính sự tồn tại ngang ngược của bức tường lại càng thổi bùng ngọn lửa khát vọng vượt ngục của người dân.\\[4pt]
\textit{(But the very brutal existence of the wall only fanned the flames of the people's desire to escape.)}

Đến tháng 11 năm 1989, sức ép của sự dồn nén đã khiến bức tường này phải sụp đổ trong một đêm không đổ máu.\\[4pt]
\textit{(By November 1989, the pressure of repression caused this wall to fall in a bloodless night.)}

Sự phủ định của bức tường cuối cùng lại trở thành nốt nhạc mở đầu cho bài ca thống nhất và dân chủ của cả châu Âu.\\[4pt]
\textit{(The negation of the wall ultimately became the opening note for the song of unification and democracy of all Europe.)}

\end{truyen}
\begin{baihoc}
  Càng bị áp bức và cấm đoán, khát vọng tự do của con người càng vươn lên mãnh liệt hơn.\\[4pt]
  \textit{(The more oppressed and forbidden, the more intensely the human thirst for freedom rises.)}
\end{baihoc}

\section{Muhammad Ali Từ Chối Đi Quân Dịch}
\begin{truyen}{Muhammad Ali Từ Chối Đi Quân Dịch}{Muhammad Ali --- Hoa Kỳ, 1967}
\chuhoa{N}{ăm} 1967, Muhammad Ali --- Đương kim Vô địch quyền anh Hạng nặng Thế giới --- nhận lệnh nhập ngũ để sang chiến đấu tại Việt Nam.\\[4pt]
\textit{(In 1967, Muhammad Ali --- the reigning World Heavyweight Boxing Champion --- received orders to be drafted to fight in Vietnam.)}

Ông kiên quyết từ chối cầm súng chống lại một dân tộc không hề làm hại mình ở cách xa nửa vòng trái đất.\\[4pt]
\textit{(He adamantly negated taking up arms against a people halfway across the globe who had done him no harm.)}

Vì tiếng ``Không'' này, Ali bị tước đai vô địch, cấm thi đấu, bị phạt tù 5 năm và mất đi những năm tháng thanh xuân chói lọi nhất của sự nghiệp.\\[4pt]
\textit{(Because of this 'No', Ali was stripped of his championship title, banned from boxing, sentenced to 5 years in prison, and lost the most prime years of his career.)}

Nhưng bằng việc phủ định một cuộc chiến tranh vô nghĩa, ông đã khẳng định bản lĩnh của một biểu tượng hòa bình và chống phân biệt phân biệt chủng tộc.\\[4pt]
\textit{(But by negating a meaningless war, he affirmed the courage of a symbol for peace and anti-racism.)}

Ali vĩ đại không chỉ vì những chiến thắng trên sàn võ, mà vì sức mạnh ông thể hiện khi làm trái lệnh các nhà cầm quyền.\\[4pt]
\textit{(Ali was great not only for his victories in the ring, but for the strength he showed in defying the commands of power.)}

\end{truyen}
\begin{baihoc}
  Kiên định với nguyên tắc đạo đức nội tâm đôi khi đòi hỏi bạn phải từ chối ngay cả vinh quang của thế gian.\\[4pt]
  \textit{(Holding fast to internal moral principles sometimes requires you to reject even the glory of the world.)}
\end{baihoc}

\section{Sự Từ Chối Của Khổng Minh Về Lương Thưởng}
\begin{truyen}{Sự Từ Chối Của Khổng Minh Về Lương Thưởng}{Gia Cát Lượng --- Trung Quốc, thế kỷ 3}
\chuhoa{G}{ia} Cát Lượng làm Thừa tướng nước Thục, quyền khuynh triều dã, nắm cả vận mệnh quốc gia trong tay.\\[4pt]
\textit{(Zhuge Liang served as Prime Minister of Shu, his power overwhelming the court, holding the nation's destiny in his hands.)}

Nhưng ông nhiều lần khước từ mọi bổng lộc, đất đai, châu báu mà Lưu Thiện ban thưởng cho ông.\\[4pt]
\textit{(But he repeatedly negated all salaries, lands, and treasures that Liu Shan bestowed upon him.)}

Ông tâu với vua: ``Thần ở Thành Đô có 800 gốc dâu, 15 khoảnh ruộng gầy, con cháu đời đời ăn mặc cũng đủ, không cần ban thêm gì.''\\[4pt]
\textit{(He reported to the king: `Your subject has 800 mulberry trees and 15 acres of poor land in Chengdu, enough for my descendants to eat and clothe for generations, I need no more.')}

Sự phủ định danh lợi và tiền tài tận cùng đó đã biến Khổng Minh thành biểu tượng bất diệt của lòng trung thành và thanh liêm.\\[4pt]
\textit{(That extreme negation of fame and wealth made Zhuge Liang an immortal symbol of loyalty and integrity.)}

Ông chứng minh rằng chính trị gia vĩ đại nhất là người không đòi hỏi bất cứ thứ gì cho bản thân.\\[4pt]
\textit{(He proved that the greatest politician is the one who demands nothing for himself.)}

\end{truyen}
\begin{baihoc}
  Người không thể bị mua chuộc bởi tư lợi cá nhân là người có sức mạnh bảo vệ lợi ích của thiên hạ.\\[4pt]
  \textit{(The one who cannot be bought by personal gain is the one with the strength to protect the public good.)}
\end{baihoc}

\section{Cuộc Biểu Tình Ngồi Của Gandhi}
\begin{truyen}{Cuộc Biểu Tình Ngồi Của Gandhi}{Mahatma Gandhi --- Ấn Độ, 1930s}
\chuhoa{K}{hác} với mọi phong trào giải phóng khác trên thế giới, Gandhi không kêu gọi người dân Ấn Độ cầm súng đánh lại người Anh.\\[4pt]
\textit{(Unlike other liberation movements in the world, Gandhi did not call on the Indian people to take up arms against the British.)}

Phương pháp của ông là 'Ahimsa' (Bất bạo động) --- tuyệt đối từ chối sử dụng bạo lực, từ chối hợp tác với chính quyền.\\[4pt]
\textit{(His method was 'Ahimsa' (Non-violence) --- absolutely negating the use of violence, negating cooperation with the authorities.)}

Người Ấn Độ bỏ việc, ngồi im trên đường phố, để mặc cho lính canh đánh đập mà không hề phản kháng.\\[4pt]
\textit{(Indians quit their jobs, sat silently on the streets, letting the guards beat them without any retaliation.)}

Sự phủ định bạo lực này lại là đòn tấn công mạnh nhất, làm tê liệt hoàn toàn đế quốc Anh và phơi bày sự tàn bạo của kẻ thống trị.\\[4pt]
\textit{(This negation of violence was the most powerful attack, completely paralyzing the British empire and exposing the brutality of the rulers.)}

Cuối cùng, không tốn một viên đạn, người Ấn Độ đã khẳng định được nền độc lập của họ trước đế quốc hùng mạnh nhất thế giới.\\[4pt]
\textit{(Ultimately, without firing a single bullet, the Indians affirmed their independence against the world's most powerful empire.)}

\end{truyen}
\begin{baihoc}
  Sức mạnh tinh thần có sức công phá mãnh liệt hơn bất kỳ vũ khí vật chất nào trên đời.\\[4pt]
  \textit{(Spiritual strength has more destructive power than any physical weapon on earth.)}
\end{baihoc}
